\documentclass{article}
\usepackage[utf8]{inputenc}
\title{Pedestrian Crossing Behavior Review}
\begin{document}
\maketitle
\section{Introduction and Research Context}
This review synthesizes evidence from 3 papers examining Pedestrian crossing behavior at unsignalized crosswalk in presence of FAV with external communication features, analyze factors affecting pedestrian red-light running behavior, predict probability of red-light running based on significant variables, with findings organized around 8 discrete claims drawn from a structured comparison matrix. The reviewed evidence employs Experimental VR simulation study, statistical modeling, random parameter logit model as primary analytical frameworks. A persistent challenge in this literature is the heterogeneity in how studies operationalize crossing behavior constructs, report metrics, and handle context specificity — factors that collectively limit direct cross-study comparison. This review structures its analysis around methodological families, task-level performance patterns, metric reporting practices, and verified gaps (1 identified), moving beyond siloed paper summaries toward a structured evidence synthesis. The review is organized as follows: Section 2 examines methodological approaches and their representational tradeoffs; Section 3 addresses task-level findings and performance patterns; Section 4 evaluates metric reporting and comparability; Section 5 identifies verified research gaps and implications; Section 6 provides a field-level synthesis and research priorities.

\section{Methodological Approaches}
The random parameter logit model family constitutes a primary methodological strand in the reviewed literature. Represents a distinct representational approach with characteristic strengths and limitations. Evidence across reviewed studies suggests that this approach is particularly well-suited to analyzing crossing decision determinants under specific traffic conditions. Limitations of this approach include reduced generalizability to contexts with substantially different traffic signal configurations. The Experimental VR simulation study family represents a primary methodological strand in the reviewed literature. Represents a distinct representational approach with characteristic strengths and limitations. Evidence across reviewed studies suggests that this approach is particularly well-suited to analyzing crossing decision determinants under specific traffic conditions. Limitations of this approach include reduced generalizability to contexts with substantially different traffic signal configurations. Across these 2 methodological families, a clear tension emerges between interpretive richness and predictive performance. Structural models — including discrete choice and social force approaches — provide parameters directly interpretable for traffic engineering and policy purposes, whereas data-driven approaches achieve higher predictive accuracy on specific benchmark datasets but lack comparable policy utility. No single paradigm has been validated across the full range of pedestrian populations, traffic environments, and signal configurations represented in this corpus.

\section{Task Scope and System Capabilities}
The reviewed literature addresses a heterogeneous set of crossing-related tasks, including analyze pedestrian red light running behavior, identify factors affecting red light violation, examine interaction effects between personal and environmental factors, analyze factors affecting pedestrian red-light running behavior, predict probability of red-light running based on significant variables. This task diversity reflects the variety of questions the field has sought to answer: from predicting individual crossing initiation decisions to evaluating system-level safety outcomes. Performance claims vary significantly across task categories. Task-level performance is most consistently reported for analyze pedestrian red light running behavior, where 8 discrete evidence claims were extracted from the reviewed studies. Reported metrics include: red light running violation binary; odds ratios; marginal effects; elasticities; AIC, Score test significance (p<0.05); odds ratios; model coefficients; confidence intervals. Metric reporting heterogeneity is substantial — some studies report normalized rates while others report absolute counts or time-based measures — making cross-study performance comparison difficult without explicit normalization. A cross-cutting observation is that task definitions themselves vary across studies: what constitutes a 'crossing violation' or an 'acceptable gap' is operationalized differently across contexts, contributing to the apparent heterogeneity in findings. Standardizing task and metric definitions would substantially improve comparability across studies.

\section{Claim Synthesis and Cross-Paper Comparability}
Reported metrics in the reviewed literature span multiple dimensions of crossing behavior and traffic system performance, including: Score test significance (p<0.05); odds ratios; model coefficients; confidence intervals, Feature ratings (1-7 Likert scale), crossing time (seconds), waiting time before crossing (seconds), pedestrian behavior questionnaire (PBQ) scores, pedestrian receptivity questionnaire for FAVs (PRQF) scores, personal innovativeness scores. Performance reporting is most complete for analyze factors affecting pedestrian red-light running behavior tasks. Metric reporting practices vary considerably across papers in two key respects. First, normalization approaches differ: some studies express outcomes as normalized rates per unit time or per pedestrian, while others report raw counts. Second, the temporal resolution of reported metrics ranges from instantaneous (e.g., crossing initiation delay at a single phase) to aggregate (e.g., violation frequency over a multi-hour observation window). These differences make direct cross-study comparison difficult without explicit standardization. Addressing this heterogeneity will require consensus on metric definitions and reporting standards within the pedestrian behavior modeling community. The absence of such standards currently forces reviewers to rely on narrative synthesis rather than formal meta-analysis, limiting the precision of conclusions that can be drawn.

\section{Research Gaps and Opportunities}
The comparative analysis of 2 papers reveals 1 persistent gaps that current approaches have not adequately resolved. These gaps are documented through verified evidence from the review matrix and are categorized by severity below. Gap [9964e351] (medium severity): There are unresolved claim differences across papers under related tasks, suggesting a need for clearer comparison criteria.. This gap limits the ability to draw robust conclusions from the existing literature. These gaps collectively indicate that the field lacks the empirical foundation needed to support confident generalization of crossing behavior models across contexts, populations, and technologies. Prioritizing the most severe gaps in future empirical work will be essential for advancing both scientific understanding and the practical applicability of crossing behavior models.

\section{Discussion: Synthesis and Future Directions}
The evidence across 2 reviewed papers converges on a field-level picture in which pedestrian crossing behavior at signalized intersections is influenced by a complex interaction of individual characteristics, traffic flow properties, signal timing parameters, and environmental context. Methodological approaches to modeling this behavior range from structural models grounded in behavioral theory to data-driven approaches optimized for predictive accuracy. Three stable conclusions emerge from this review. First, waiting time is among the most consistent predictors of crossing initiation and violation behavior across reviewed studies. Second, methodological approaches exhibit a persistent tradeoff between interpretability and predictive performance. Third, metric reporting heterogeneity across studies substantially limits the precision of cross-study synthesis. Two persistent tensions run through this literature. The first concerns the cross-context validity of gap acceptance thresholds: values derived from studies in North American and European contexts may not transfer to other traffic environments with different vehicle-pedestrian interaction norms. The second tension concerns the behavioral effects of automated vehicles: whether pedestrians recalibrate their crossing heuristics with sustained AV exposure remains unresolved due to the absence of longitudinal field studies. Resolving these tensions will require longitudinal field studies in diverse cultural contexts. Priority should be given to addressing Gap [0] concerning cross-cultural threshold transfer and Gap [1] concerning longitudinal behavioral adaptation in AV environments — gaps that carry direct implications for both road safety policy and the design of pedestrian-accepting automated vehicle systems.

\section{Conclusion}
The evidence across 2 reviewed papers converges on a field-level picture in which pedestrian crossing behavior at signalized intersections is influenced by a complex interaction of individual characteristics, traffic flow properties, signal timing parameters, and environmental context. Methodological approaches to modeling this behavior range from structural models grounded in behavioral theory to data-driven approaches optimized for predictive accuracy. Three stable conclusions emerge from this review. First, waiting time is among the most consistent predictors of crossing initiation and violation behavior across reviewed studies. Second, methodological approaches exhibit a persistent tradeoff between interpretability and predictive performance. Third, metric reporting heterogeneity across studies substantially limits the precision of cross-study synthesis. Two persistent tensions run through this literature. The first concerns the cross-context validity of gap acceptance thresholds: values derived from studies in North American and European contexts may not transfer to other traffic environments with different vehicle-pedestrian interaction norms. The second tension concerns the behavioral effects of automated vehicles: whether pedestrians recalibrate their crossing heuristics with sustained AV exposure remains unresolved due to the absence of longitudinal field studies. Resolving these tensions will require longitudinal field studies in diverse cultural contexts. Priority should be given to addressing Gap [0] concerning cross-cultural threshold transfer and Gap [1] concerning longitudinal behavioral adaptation in AV environments — gaps that carry direct implications for both road safety policy and the design of pedestrian-accepting automated vehicle systems.

\end{document}
% BIB START
@article{05_2018_deb_trf,
  author={Shuchisnigdha Deb; Lesley J. Strawderman; Daniel W. Carruth},
  title={Investigating pedestrian suggestions for external features on fully autonomous vehicles: A virtual reality experiment},
  year={2018},
  journal={Transportation Research Part F}
}

@article{12_2016_zhang_aap,
  author={Weihua Zhang; Kun Wang; Lei Wang; Zhongxiang Feng; Yingjie Du},
  title={Exploring factors affecting pedestrians' red-light running behaviors at intersections in China},
  year={2016},
  journal={Accident Analysis and Prevention}
}

@article{13_2020_chen_trf,
  author={Dianchen Zhu; N.N. Sze; Lu Bai},
  title={Roles of personal and environmental factors in the red light running propensity of pedestrian: Case study at the urban crosswalks},
  year={2020},
  journal={Transportation Research Part F}
}